\section{Conclusão e Trabalhos Futuros}
\label{sec:conclusao}


Este trabalho apresentou uma abordagem algorítmica e topológica, baseada em grafos heterogêneos multicamada, para a identificação de indícios de conluio e simulação de competitividade em contratações públicas municipais. A aplicação aos dados abertos do SICOM/TCE-MG, referentes ao exercício de 2025 do município de Divinópolis/MG, demonstrou o potencial da integração entre registros licitatórios, informações cadastrais da Receita Federal, obtidas por meio da BrasilAPI, e técnicas de análise estrutural para identificar relações que não são evidentes em consultas isoladas às bases de dados.

As principais contribuições do estudo podem ser sintetizadas em três aspectos.

\subsection{Eficiência da Redução Topológica}

A combinação entre a filtragem por escore multidimensional ($\theta \geq 2{,}00$), a remoção de vértices isolados e a identificação de componentes conexas permitiu reduzir a rede original, composta por 1.788 vértices e 5.914 arestas, a um subgrafo de interesse pericial com apenas 26 vértices e 23 arestas. Essa transformação corresponde a uma redução de 98,55\% no número de vértices e de 99,61\% no número de arestas.

O procedimento permitiu eliminar elementos sem participação nas relações selecionadas, reduzindo a interferência da estrutura geral da rede e concentrando a análise em seis componentes conexas prioritárias para auditoria. Dessa forma, a metodologia demonstrou que a redução estrutural determinística pode tornar redes de contratações públicas mais interpretáveis, preservando as relações consideradas relevantes segundo os critérios de pontuação estabelecidos.

\subsection{Identificação de Indícios de Controle e Vínculos Compartilhados}

A incorporação de informações do Quadro de Sócios e Administradores (QSA), combinada à análise de relações societárias e à identificação de motivos estruturais, permitiu localizar três motivos triangulares na representação utilizada para a análise, distribuídos em dois núcleos críticos. Entre os casos identificados, destaca-se a ocorrência de relações societárias compartilhadas envolvendo sócio pessoa física e uma holding controladora, associadas à disputa simultânea de 59 itens em um mesmo certame.

Além das coincidências societárias, a ponderação onomástica inspirada no princípio IDF e a normalização de endereços e telefones permitiram identificar quatro núcleos adicionais sem coincidência formal de sócios no QSA. Esses resultados evidenciam a utilidade da combinação de diferentes atributos cadastrais e relacionais para a identificação de possíveis vínculos domiciliares, operacionais e familiares que poderiam não ser detectados por cruzamentos cadastrais convencionais.

\subsection{Contribuição ao Controle Externo}

O pipeline desenvolvido constitui uma ferramenta de apoio à triagem de procedimentos licitatórios, oferecendo critérios explícitos de pontuação, relações rastreáveis e resultados passíveis de inspeção. Diferentemente de abordagens baseadas exclusivamente em modelos opacos, a metodologia permite relacionar os indícios identificados aos atributos e às conexões que contribuíram para sua classificação.

Ressalta-se, contudo, que os resultados possuem caráter indiciário. A identificação de relações societárias, cadastrais ou estruturais não constitui, isoladamente, comprovação de conluio ou de irregularidade. Os resultados devem, portanto, orientar procedimentos posteriores de auditoria e investigação, com verificação documental e análise contextual pelos órgãos competentes. A utilização de identificadores pseudonimizados também foi considerada na arquitetura, em alinhamento com os princípios de proteção de dados estabelecidos pela Lei Geral de Proteção de Dados Pessoais (Lei nº 13.709/2018).

\subsection{Trabalhos Futuros}

Como continuidade desta pesquisa, identificam-se três direções principais. A primeira consiste na incorporação de atributos temporais e financeiros das propostas comerciais e dos lances registrados nas atas de julgamento, especialmente nos dados de \texttt{julgLicitacao.csv}. Essa extensão permitiria investigar padrões de alternância de vencedores (\textit{bid rotation}) e apresentação de propostas de cobertura (\textit{cover bidding}), ampliando a análise para além das relações cadastrais e societárias.

A segunda direção consiste na integração de métricas de centralidade e algoritmos de detecção de comunidades em bases plurianuais do SICOM. Essa abordagem possibilitaria acompanhar a evolução temporal dos grupos econômicos, identificar mudanças nas relações entre fornecedores e órgãos contratantes e investigar a persistência de padrões suspeitos ao longo de diferentes exercícios financeiros.

Por fim, pretende-se investigar a aplicação de modelos de aprendizado de representação em grafos heterogêneos (\textit{Heterogeneous Graph Neural Networks}), com o objetivo de explorar mecanismos de calibração dos pesos do escore multidimensional a partir de conjuntos históricos de casos analisados ou julgados pelos órgãos de controle. Essa possibilidade deverá ser avaliada considerando a disponibilidade de dados rotulados, a interpretabilidade dos modelos e a necessidade de preservar a rastreabilidade dos critérios utilizados na identificação dos indícios.
